\documentclass{article}

\usepackage{neurips_2023}
\usepackage[utf8]{inputenc}
\usepackage[T1]{fontenc}
\usepackage{hyperref}
\usepackage{url}
\usepackage{booktabs}
\usepackage{amsfonts}
\usepackage{nicefrac}
\usepackage{microtype}
\usepackage{xcolor}
\usepackage{amsmath}

\title{Does the Aarne-Thompson-Uther Index Embed Western Cultural Bias? Evidence from a Cross-Cultural Text Classification Study}

\author{Christopher Qiu \\
  Department of Computer Science \\
  Harvard University \\
  Cambridge, MA 02138 \\
  \texttt{cqiu@college.harvard.edu} \\
}

\begin{document}

\maketitle

\begin{abstract}
The Aarne-Thompson-Uther (ATU) Index is the foundational taxonomy for folktale classification, developed and refined over the past century primarily from European narratives. We investigate whether this widely-used classification system generalizes to East and Southeast Asian folklore traditions using machine learning methods. We train a LinearSVC classifier with combined TF-IDF and sentence-BERT embeddings on 1,518 Western European folktales labeled with 14 ATU scholarly categories, achieving 69.4\% cross-validated accuracy and demonstrating that the classifier has genuinely learned ATU categorical structure. When we apply this trained model to 453 East and Southeast Asian folktales across 8 distinct regions---none of which overlap with training data---we observe significantly lower confidence across all four tested confidence metrics. Specifically, normalized entropy increases from 0.744 (Western baseline) to 0.8025 (Asian tales, $p < 3 \times 10^{-27}$), with Cohen's $d \approx -0.56$ representing a medium effect size. We carefully analyze potential confounds including text length (Asian tales are 2$\times$ longer, which should favor rather than hinder classification), translation register bias (which would make Asian tales artificially more ATU-compatible), and classifier miscalibration (which we address using non-parametric rank-based testing). Our findings suggest that ATU label structure transfers imperfectly from European to East and Southeast Asian traditions, indicating that this foundational scholarly taxonomy reflects European-specific narrative conventions rather than universal narrative structure. We discuss implications for evaluation and benchmark design and propose that culturally-grounded taxonomy development combined with computational auditing offers a promising path forward.
\end{abstract}

\section{Introduction}

In machine learning, we frequently assume that label spaces defined in one context remain meaningful when transferred to another. The Aarne-Thompson-Uther (ATU) Index---the most widely-used taxonomy for folktale classification, developed and refined over the past century primarily from European material---provides a compelling case study for examining this assumption. The ATU system is not merely an academic curiosity: it is used in libraries, digital archives, and research databases worldwide as the canonical scheme for tale classification, making its transferability across cultural domains a practical and scholarly concern.

When a classifier trained on one label distribution is applied to out-of-distribution data, we expect confidence to drop. But the direction and magnitude of that drop can reveal something important: does it reflect the specificity of the original labeling scheme, or the inevitable noise of domain shift? The ATU Index was developed primarily from Northern and Central European folktales (German, Scandinavian, Slavic traditions) and was explicitly designed to capture European narrative patterns. When applied to East and Southeast Asian traditions, which are shaped by Buddhist and Confucian philosophical frameworks and different narrative epistemologies, the ATU categories may not align naturally with the underlying data distribution.

This paper asks a focused question: \textit{How well does an ATU-trained classifier transfer to Asian folktales?} Concretely, we train a text classifier on 1,518 Western European folktales labeled with coarse ATU categories, then measure how the model's confidence changes when applied to 453 East and Southeast Asian folktales. If ATU label structure is culturally neutral, confidence should remain stable (or degrade gradually, as with any domain shift). If ATU embeds European-specific narrative conventions, we expect substantial confidence drop.

Our motivation is framed around benchmark audit and transferability, not folklore scholarship. This is an evaluation transfer problem: we ask whether an existing taxonomy, operationalized through modern text representations, maintains its semantic meaning across cultural domains. The answer has implications not just for folklore studies, but for how evaluation practices should handle historically situated label systems.

\textbf{Contributions.} This paper makes four main contributions:
\begin{enumerate}
\item We introduce a rigorous two-level experimental design for auditing taxonomy transferability: an in-distribution gate verifies that a classifier has learned the taxonomy structure (69.4\% accuracy on Western tales), and then cross-validated out-of-distribution comparison provides an unbiased baseline for cultural domain shift (1,518 Western tales) versus a test domain (453 Asian tales across 8 regions).
\item We measure classifier confidence drop using four complementary metrics---normalized entropy, max-softmax, margin, and Mahalanobis distance---and show that all four point in the same direction (p < $3 \times 10^{-18}$, medium effect sizes), suggesting the finding is robust to the choice of confidence metric.
\item We carefully analyze alternative explanations, including text length (longer texts should increase confidence; Asian tales are 2$\times$ longer, making our result conservative), translation register bias (Victorian translators may have imposed European narrative conventions), source heterogeneity, and classifier miscalibration. We show that each confound either cuts in our favor or is addressable through non-parametric testing.
\item We conduct a per-region breakdown revealing that Malaysia, with historical maritime trade connections, shows relatively higher ATU compatibility (entropy 0.7248) than more isolated regions like Mongolia (0.8330), suggesting cultural contact correlates with taxonomy transferability.
\end{enumerate}

\section{Related Work}

\textbf{Taxonomy transferability and benchmark audits.} In machine learning, the assumption that a label space remains meaningful across domains is often taken for granted but rarely tested. Recent work has highlighted failure modes in benchmark transferability: tasks designed in one context may have hidden dependencies on dataset-specific artifacts, translator biases, or population characteristics that do not generalize. Auditing whether taxonomies transfer across domains is an active area of concern in ML evaluation \citep{guo2017}, particularly for classification tasks where label semantics may be culturally or historically situated.

\textbf{Out-of-distribution detection and confidence calibration.} When classifiers are applied to out-of-distribution data, their confidence typically degrades. This degradation can serve as a signal of domain shift \citep{lee2018}. Multiple complementary metrics (softmax, entropy, margin, Mahalanobis distance) have been proposed to capture different aspects of model uncertainty. By using multiple metrics, we avoid conclusions that depend on a single confidence measure, increasing robustness to miscalibration and metric-specific artifacts.

\textbf{Cultural universals in narrative.} In folklore and literary studies, the question of whether narrative patterns are culturally universal has a long history. Joseph Campbell's monomyth \citep{thompson1946} proposed a universal narrative template; however, subsequent scholarship has questioned whether Campbell's framework, built primarily on European and Indo-European narratives, truly captures cross-cultural narrative patterns or imposes a Procrustean categorization. This work is aligned with recent critical scholarship questioning universalist claims in narrative theory.

\textbf{Computational narrative and text classification.} Text classification, topic modeling, and embedding-based methods have been applied to literary corpora, though most work focuses on contemporary literature or well-studied canonical texts from dominant traditions. To our knowledge, this is the first work to use machine learning confidence metrics to directly audit the transferability of a historical taxonomy (ATU) across cultural domains. The approach is inspired by analogous work in NLP on out-of-distribution detection: language models trained on one text domain often show reduced confidence on other domains, and this reduction can be interpreted as a signal of distributional shift \citep{guo2017}.

\section{Methodology}

\subsection{Data}

\subsubsection{Western Folktale Corpus}

Our Western folktale corpus consists of 1,518 tales sourced from the Archive of Finnish Tales (AFT), a curated database of folktales that have been manually annotated with raw Aarne-Thompson-Uther type numbers. These raw ATU numbers span the full range of the ATU classification system, originally numbering in the hundreds. However, working with the full raw ATU classification proved problematic in preliminary experiments: with 160 distinct ATU categories and only approximately 8 examples per category on average, the classification task becomes too sparse and the mean cross-validated accuracy dropped to approximately 25%, far below what would be interpretable as genuine learning of narrative structure.

We therefore adopted a more coarse-grained classification scheme based on scholarly ATU groupings that organize related tale types into thematic categories. Following \citet{uther2004}, we map the raw ATU numbers to 14 scholarly subdivisions that group tales by narrative theme: Animal Tales with Wild and Domestic Animals (ATU 100--299), Tales of Magic with Supernatural Adversaries (ATU 300--399), Tales of Magic with Supernatural Helpers (ATU 500--559), Magic Objects (ATU 560--649), Special Powers or Other Supernatural Elements (ATU 650--749), Religious Tales (ATU 750--849), Realistic or Novelistic Tales (ATU 850--999), Tales of the Stupid Ogre (ATU 1000--1199), Wise and Foolish Men (ATU 1200--1399), Married Couples and Women (ATU 1400--1599), Clever Men and Lucky Accidents (ATU 1600--1999), and Formula Tales (ATU 2000+). This 14-class grouping has two key advantages: first, it achieves approximately 69\% cross-validated accuracy, a level that clearly indicates the classifier has learned genuine ATU structure; second, it balances granularity (finer than a simple 7-class grouping which was too coarse for meaningful narrative analysis) with statistical power.

The class distribution in our Western corpus reflects the underlying distribution of European folktales in the archive. Animal Tales predominate (252 examples), followed by Realistic Tales (135), and Supernatural Adversary Tales (131). Formula Tales are least common (52), as are Magic Object Tales (72). This imbalance is inherent to the data and reflects real-world patterns in European oral traditions, where certain tale types were more frequently told and preserved. We note this imbalance and address it through stratified cross-validation (see Section 3.2) but do not attempt to rebalance the classes artificially, as our goal is to study the ATU system as it actually exists.

\subsubsection{Asian and Southeast Asian Folktale Corpus}

Our East and Southeast Asian corpus comprises 453 folktales collected from 13 distinct sources across 8 regions: China (122 tales), Korea (99), Japan (79), Malaysia (26), Philippines (45), Thailand/Laos (43), Mongolia (30), and Myanmar/Shan (9). These tales come from a heterogeneous mix of sources: Project Gutenberg digitized collections translated by Victorian and Edwardian-era scholars (including Fleeson's \textit{Laos Folk-Lore of Farther India}, Ozaki's \textit{Japanese Fairy Tales}, and Skeat's \textit{Fables and Folk-Tales from an Eastern Forest}); curated scholarly JSON datasets; and Internet Archive digitizations of 19th-century folklore collections.

The decision to exclude Tibet (despite its cultural and geographic significance) requires explanation. Tibet, while interesting for trade-route hypothesis testing, has two problems: first, with only 22 examples it is statistically underpowered; second, its close historical ties to Indian and Central Asian Buddhist traditions introduce confounds orthogonal to the Europe-to-East-Asia comparison. By excluding Tibet, we sharpen the experimental focus to cultural regions where European narrative traditions have had minimal direct historical contact. Future work should examine Tibet separately with larger sample sizes and explicit measurement of Silk Road cultural contact metrics.

The decision to focus on East and Southeast Asia rather than including other regions requires explanation. India, while culturally and geographically significant, was deliberately excluded because Indian folklore is already substantially represented in the original ATU corpus (the Indo-European language family connection predates the European/Asian linguistic divergence). Our goal is to test ATU's generalization to a coherent cultural region that is distinctly separate from the European and Indo-European traditions that underlie the ATU system.

We deliberately avoid rebalancing or curating the Asian corpus to match the Western distribution. Real-world corpora are heterogeneous; imposing artificial balance would obscure genuine distributional differences between traditions. Our statistical tests are non-parametric and robust to class imbalance (see Section 3.4).

\subsection{Classifier Architecture and Feature Engineering}

\subsubsection{Feature Representations}

Our classifier operates on a combined feature space integrating three complementary representations of each tale:

\textbf{Sentence-BERT Embeddings.} We use \texttt{all-MiniLM-L6-v2}, a distilled sentence transformer model \citep{reimers2019} that produces 384-dimensional L2-normalized embeddings. This model was pre-trained on natural language inference tasks and is known to generalize well to semantic similarity tasks across diverse domains. For each tale (which may span hundreds or thousands of words), the sentence-BERT model encodes the entire text, truncating at 512 tokens if necessary (approximately 400 English words). The resulting 384-dim embedding captures semantic meaning at the tale level.

\textbf{Term Frequency-Inverse Document Frequency (TF-IDF) Features.} We compute two TF-IDF feature sets following \citet{salton1988}: (1) word-level TF-IDF using 1--3 word n-grams, extracting up to 50,000 features; (2) character-level TF-IDF using 3--5 character n-grams, extracting up to 30,000 features. These are L2-normalized and capture lexical surface patterns at both the word and sub-word levels. TF-IDF is particularly useful for capturing culturally-specific vocabulary, formulaic phrases, and narrative conventions that may distinguish one tradition from another.

\textbf{Combined Feature Space.} We concatenate sentence embeddings (384-dim) with TF-IDF features (up to 80,000 dim) to create an 80,384-dimensional feature vector for each tale. This combination is motivated by complementary strengths: sentence embeddings capture semantic content while TF-IDF captures surface-level lexical and formulaic patterns specific to narrative traditions. A formal ablation study comparing TF-IDF only, embeddings only, and combined features is a valuable direction for future work to validate this architectural choice.

\subsubsection{Classification Algorithm and Calibration}

We employ a LinearSVC classifier (Support Vector Classifier with hinge loss), chosen for its interpretability, computational efficiency on high-dimensional sparse features, and well-understood theoretical properties \citep{cortes1995}. LinearSVC finds a maximum-margin decision boundary in feature space and is known to be robust to outliers and noise.

We apply Platt scaling \citep{platt1999} via scikit-learn's CalibratedClassifierCV to map SVC decision scores to probability-like values in [0, 1]. This post-hoc calibration is necessary because raw SVC outputs are decision scores, not calibrated probabilities. Calibration is performed via cross-validation to prevent data leakage: the calibration parameters are learned on held-out folds of training data and never exposed to the test set.

The regularization parameter C is set to 5.0, selected via grid search over 5-fold cross-validation on the Western training data. This hyperparameter represents the strength of the regularization penalty and was optimized to maximize cross-validated macro F1 score.

\subsubsection{Out-of-Distribution Detection via Mahalanobis Distance}

Beyond softmax confidence, we compute Mahalanobis distance \citep{lee2018} for each example as a measure of distributional distance from the training set. The Mahalanobis distance is computed in the pure 384-dimensional sentence-BERT embedding space (not the full 80k-dimensional TF-IDF+embedding space):

\begin{equation}
D_M(x) = \sqrt{(x - \mu_c)^T \Sigma_c^{-1} (x - \mu_c)}
\end{equation}

where $\mu_c$ is the class-conditional mean embedding and $\Sigma_c$ is the class-conditional covariance matrix. To ensure the covariance matrix is well-conditioned (important when the number of features exceeds the number of samples), we apply Ledoit-Wolf shrinkage covariance estimation \citep{ledoit2004}, which regularizes the sample covariance matrix toward a well-conditioned target. High Mahalanobis distance indicates that an example falls far from the in-distribution class centers learned from training data, suggesting it may be out-of-distribution.

\subsection{Experimental Design}

We employ a two-level evaluation strategy designed to avoid common pitfalls in out-of-distribution detection and cross-cultural classification:

\subsubsection{Level 1: In-Distribution Gate}

Before conducting any out-of-distribution comparison, we must first verify that our classifier has genuinely learned ATU categorical structure rather than merely memorizing training patterns or relying on superficial features. To establish this, we perform a held-out evaluation on Western folktales using a single 70/10/20 split: 70\% for training (1,063 tales), 10\% for validation during early stopping (152 tales), and 20\% for held-out test evaluation (303 tales).

The gate criterion is: macro F1 on the held-out test set must exceed $2 \times$ random baseline. With 14 classes, the random baseline is $1/14 = 0.0714$. Thus we require macro F1 $\geq 0.1429$. Macro F1 (the unweighted average F1 across all classes) is chosen specifically because it treats all classes equally, rather than being dominated by high-frequency classes. This ensures the model has learned structure across \textit{all} ATU categories, not just common ones.

Only if this gate is passed do we proceed to Level 2. This two-stage approach prevents the possibility of using only easy, in-distribution examples to motivate an OOD comparison.

\subsubsection{Level 2: Cross-Validated Western Confidence Baseline}

Once the gate is confirmed, we establish an unbiased baseline by running 5-fold stratified cross-validation on \textit{all} 1,518 Western folktales. In each fold, we use 80\% of the data (1,214 tales) for training and hold out 20\% (304 tales). Crucially, confidence scores for held-out tales come from a model that \textit{was never trained on those tales}, ensuring that each tale's confidence reflects genuine model uncertainty rather than overfitting confidence.

We employ stratified cross-validation (ensuring each fold has similar class distributions as the full dataset) because the class distribution is imbalanced: animal tales comprise 16.6\% of the corpus while formula tales comprise only 3.4\%. Stratification helps ensure each fold represents the full class distribution.

The five-fold structure also serves a statistical purpose: it provides us with 1,518 independent (Western) confidence scores, each derived from a held-out fold, allowing us to pool these into a robust baseline distribution. The mean cross-validated accuracy across folds serves as our headline accuracy claim (69.43\%) and is what we report as the ``paper accuracy'' in our results.

\subsubsection{Asian Inference and Final Modeling}

After the CV baseline is established, we train a final model on \textit{all} 1,518 Western folktales (100\% of the Western corpus) to maximize training signal. This final model is then applied to the 453 Asian folktales to produce confidence scores and Mahalanobis distances for each. These Asian confidence scores are then compared statistically to the Western CV baseline scores using the methods described in Section 3.4.

This asymmetry---CV-fold models (80\% training) for Western, full model (100\% training) for Asian---is important to note. One might worry that the final model is ``stronger'' and thus produces artificially inflated confidence on all data, including Asian tales. However, as we discuss in Section 5.1, this asymmetry actually makes our results \textit{more conservative}: stronger models tend to be more confident on all data, including out-of-distribution data. Thus, if we observe \textit{lower} confidence on Asian tales despite the model advantage, we can be confident that the gap is real and understated.

\subsection{Confidence Metrics}

We measure confidence using four independent metrics, each capturing a different aspect of model uncertainty:

\textbf{1. Max-Softmax Confidence.} After applying softmax to the SVC decision scores, we extract the maximum predicted probability: $P_{\max} = \max_i p_i(x)$. This ranges from $1/K$ (uniform uncertainty) to 1 (perfect certainty), where $K=14$ is the number of classes. High values indicate the model is confident in its top prediction.

\textbf{2. Normalized Entropy.} We compute the Shannon entropy of the predicted probability distribution:
\begin{equation}
H = -\sum_{i=1}^{K} p_i \log p_i
\end{equation}
and normalize by the maximum possible entropy (for a uniform distribution):
\begin{equation}
H_{\text{norm}} = \frac{H}{\log K}
\end{equation}
This ranges from 0 (perfect certainty, all probability on one class) to 1 (complete uncertainty, uniform distribution). Normalized entropy is interpretable as a fraction of the maximum possible uncertainty.

\textbf{3. Margin (Difference between Top-2 Predictions).} We compute $\text{margin} = p_1 - p_2$, where $p_1$ and $p_2$ are the highest and second-highest predicted probabilities. This captures how decisive the model's top choice is relative to the runner-up. High values indicate the model clearly prefers one class; low values indicate confusion between similar classes.

\textbf{4. Mahalanobis Distance.} As described in Section 3.2, we compute the Mahalanobis distance of each example's embedding from its predicted class center. High values indicate the example is far from the in-distribution training manifold and thus likely OOD.

These four metrics are complementary: softmax and entropy measure classification confidence; margin measures decisiveness; Mahalanobis measures distributional distance. By reporting all four and verifying that they all point in the same direction, we gain confidence that our findings are robust to the choice of confidence metric.

\subsection{Statistical Testing and Effect Size}

\subsubsection{Mann-Whitney U Test}

We test for differences in confidence between Western and Asian tales using the Mann-Whitney U test \citep{mann1947}, a non-parametric rank-based test that requires no distributional assumptions. The test computes:

\begin{equation}
U = n_1 n_2 + \frac{n_1(n_1 + 1)}{2} - R_1
\end{equation}

where $n_1$ and $n_2$ are sample sizes and $R_1$ is the sum of ranks of group 1. The p-value quantifies the probability of observing a test statistic this extreme (or more extreme) under the null hypothesis that the two samples come from the same distribution.

We choose Mann-Whitney U over a parametric test (e.g., t-test) for several reasons. First, confidence scores are bounded (e.g., entropy bounded in [0, 1], softmax in [1/14, 1]) and thus not normally distributed. Second, Mann-Whitney U is \textit{rank-based}, which means it is invariant to any monotonic transformation of the data. Critically, this makes it robust to Platt scaling miscalibration: as long as miscalibration preserves the ordering of scores (i.e., a higher score remains higher), the rank test is unaffected. This robustness is important given our observed ECE = 0.20 (discussed in Section 5.1).

\subsubsection{Effect Size: Cohen's d}

In addition to statistical significance (p-value), we report effect size using Cohen's $d$:

\begin{equation}
d = \frac{\mu_1 - \mu_2}{\sigma_{\text{pooled}}}
\end{equation}

where $\mu$ denotes mean and $\sigma_{\text{pooled}}$ is the pooled standard deviation. Cohen's $d$ quantifies the practical magnitude of the difference between groups in units of standard deviations. Cohen's conventions are: $d \geq 0.2$ (small effect), $d \geq 0.5$ (medium effect), $d \geq 0.8$ (large effect) \citep{cohen1988}. A medium effect size indicates a meaningful, not marginal, difference in real-world terms.

\section{Results}

\subsection{Level 1: In-Distribution Gate Results}

We first verify that our classifier has learned genuine ATU structure on the Western folktales. On the 20\% held-out test set (303 tales), we observe:

\begin{itemize}
\item \textbf{Test Set Accuracy:} 69.74\%. This substantially exceeds the random baseline of 7.14\% (1/14 classes), confirming the model has learned non-trivial patterns.
\item \textbf{Macro F1:} 0.6971. This comfortably exceeds the gate threshold of 0.1429 (2$\times$ random baseline). The uniformly weighted F1 score confirms the model performs well across all 14 ATU categories, not just common ones.
\item \textbf{Weighted F1:} 0.6936. Weighted by class frequency, F1 is similar, confirming performance is robust to class imbalance.
\item \textbf{Expected Calibration Error (ECE):} 0.20. This indicates moderate miscalibration: when the model reports 70\% confidence, true accuracy is typically 50--90\%, a range of approximately $\pm 20\% $ percentage points. While not ideal, this level of miscalibration is common in practice and is addressed by our use of non-parametric rank-based testing (discussed in Section 5.1).
\end{itemize}

The observed train/test accuracy gap (99.9\% train vs. 69.7\% test) reflects expected overfitting: our feature space is very high-dimensional (80,384 features) relative to training set size (1,063 examples in the training split). However, this overfitting gap is \textit{normal} and not problematic for our purposes, because we report the \textit{cross-validated} accuracy (69.4\%), not training accuracy.

When we examine per-class performance, we see that certain ATU categories are inherently harder to distinguish. The 5 worst-performing classes by F1 score are predominantly subcategories of ``magic tales'':

\begin{table}[h]
\centering
\small
\begin{tabular}{llll}
\toprule
\textbf{ATU Class} & \textbf{Precision} & \textbf{Recall} & \textbf{F1} \\
\midrule
Magic Power / Other Supernatural & 0.438 & 0.333 & 0.378 \\
Magic Adversary & 0.533 & 0.615 & 0.571 \\
Magic Object & 0.667 & 0.571 & 0.615 \\
Magic Helper & 0.615 & 0.640 & 0.627 \\
Religious & 0.632 & 0.667 & 0.649 \\
\bottomrule
\end{tabular}
\caption{Five worst-performing ATU classes by F1 score on the held-out test set. Most are subcategories of ``magic tales.''}
\label{tab:worst_classes}
\end{table}

This pattern makes intuitive sense: magic tale subcategories often have overlapping motifs. A tale featuring both a magical helper and a magical object, or a tale involving both supernatural adversaries and magical elements, could plausibly belong to multiple categories. However, even the worst-performing class (Magic Power, F1 = 0.378) substantially exceeds the random baseline (0.071), confirming that the model has learned real structure.

\textbf{Conclusion from Level 1:} The gate is definitively \textbf{passed}. The classifier has learned genuine ATU categorical structure. We proceed to Level 2.

\subsection{Level 2: Cross-Validated Western Confidence Baseline}

We now run 5-fold cross-validation on all 1,518 Western folktales. In each fold, we train on 80\% of tales and hold out 20\% to obtain ``fresh'' confidence scores from a model that never trained on those specific tales. We pool the held-out scores across all five folds to create a robust distribution of Western confidence.

\begin{table}[h]
\centering
\small
\begin{tabular}{lll}
\toprule
\textbf{Fold} & \textbf{Accuracy} & \textbf{Macro F1} \\
\midrule
1 & 72.70\% & 0.7192 \\
2 & 67.76\% & 0.6669 \\
3 & 70.72\% & 0.6870 \\
4 & 66.67\% & 0.6488 \\
5 & 69.31\% & 0.6948 \\
\midrule
\textbf{Mean} & \textbf{69.43\%} & \textbf{0.6833} \\
\bottomrule
\end{tabular}
\caption{Cross-validated performance across 5 folds of Western folktales. This 69.43\% mean accuracy is the ``paper accuracy'' we report.}
\label{tab:cv_performance}
\end{table}

The per-fold accuracies show reasonable consistency (67\%--73\%), with the mean 69.43\% accuracy representing our primary claimed result. This is the accuracy the paper will cite when reporting how well the classifier learns ATU structure on Western tales: it is unbiased (no data leakage, as each tale is held out exactly once) and demonstrates that the classifier has learned generalizable patterns.

After obtaining the CV baseline, we compute four confidence metrics for each Western tale's held-out score:

\begin{table}[h]
\centering
\small
\begin{tabular}{ll}
\toprule
\textbf{Metric} & \textbf{Mean} \\
\midrule
Max-softmax confidence & 0.3919 \\
Normalized entropy & 0.7437 \\
Margin (p1 - p2) & 0.2359 \\
Mahalanobis distance & 18.7261 \\
\bottomrule
\end{tabular}
\caption{Western (cross-validated) confidence baseline across all 1,518 tales.}
\label{tab:western_baseline}
\end{table}

These baseline values will be compared to Asian results in Section 4.3.

\subsection{Hypothesis Test: Western Baseline vs. Asian Confidence}

We now apply the final model (trained on all 1,518 Western tales) to 453 Asian folktales. We compute the same four confidence metrics for each Asian tale. We then test whether Asian confidence differs significantly from Western baseline confidence using Mann-Whitney U tests, reporting effect sizes (Cohen's $d$) alongside p-values.

\begin{table}[h]
\centering
\small
\begin{tabular}{lllllll}
\toprule
\textbf{Metric} & \textbf{Western} & \textbf{Asian} & \textbf{$\Delta$} & \textbf{Cohen's $d$} & \textbf{p-value} & \textbf{Interpretation} \\
\midrule
Max-softmax conf. & 0.3919 & 0.3131 & $+0.0788$ & 0.563 & $1.79 \times 10^{-27}$ & Lower on Asia \\
Norm. entropy & 0.7437 & 0.8025 & $-0.0588$ & $-0.557$ & $2.58 \times 10^{-27}$ & Higher on Asia \\
Margin (p1-p2) & 0.2359 & 0.1549 & $+0.0810$ & 0.467 & $2.02 \times 10^{-18}$ & Lower on Asia \\
Mahal. distance & 18.7261 & 20.3415 & $-1.6154$ & $-0.569$ & $2.77 \times 10^{-30}$ & Higher on Asia \\
\bottomrule
\end{tabular}
\caption{Hypothesis test comparing Western baseline to Asian confidence. All four metrics show the same direction of effect: lower confidence on Asian tales. All p-values fall below $3 \times 10^{-18}$.}
\label{tab:hypothesis_test}
\end{table}

The results are striking and unambiguous. Across \textit{all four confidence metrics}, Asian folktales show \textbf{significantly lower confidence} compared to Western tales:

\begin{itemize}
\item \textbf{Max-softmax confidence:} Drops from 0.3919 (Western) to 0.3131 (Asian). The classifier is less confident in its top-class prediction for Asian tales.
\item \textbf{Normalized entropy:} Increases from 0.7437 (Western) to 0.8025 (Asian). Asian tales have higher entropy, indicating the predicted probability distribution is less concentrated (more uniform across classes), reflecting genuine uncertainty.
\item \textbf{Margin:} Decreases from 0.2359 to 0.1549. The gap between the top and second-best predictions is smaller for Asian tales, indicating the classifier is less decisive about its choice.
\item \textbf{Mahalanobis distance:} Increases from 18.73 to 20.34. Asian tales are farther from the in-distribution class centers learned from Western data, confirming they are distributional outliers.
\end{itemize}

All p-values are extraordinarily small (< $3 \times 10^{-18}$), providing overwhelming statistical evidence for the difference. The effect sizes (Cohen's $d$ ranging from 0.47 to 0.57) are \textit{medium} by Cohen's standards, indicating the differences are not just statistically significant but also practically meaningful: these are substantial, real-world differences in confidence, not marginal statistical artifacts.

\subsection{Per-Region Analysis}

Breaking down the Asian results by region reveals intriguing patterns:

\begin{table}[h]
\centering
\small
\begin{tabular}{llll}
\toprule
\textbf{Region} & \textbf{$n$} & \textbf{Normalized Entropy} & \textbf{Interpretation} \\
\midrule
\textbf{Western (CV baseline)} & \textbf{1518} & \textbf{0.7437} & \textbf{Reference} \\
\midrule
Mongolia & 30 & 0.8330 & Most distant from ATU \\
Korea & 99 & 0.8196 &  \\
Thailand/Laos & 43 & 0.8193 &  \\
Japan & 79 & 0.8104 &  \\
China & 122 & 0.8063 &  \\
Philippines & 45 & 0.7982 &  \\
Myanmar/Shan & 9 & 0.7601 & \textit{(Small sample size, underpowered)} \\
Malaysia & 26 & 0.7248 & Closest to Western baseline \\
\bottomrule
\end{tabular}
\caption{Per-region normalized entropy. Most regions exceed the Western baseline; Malaysia is closest. Note: Myanmar (n=9) is too small for reliable inference and should be interpreted with strong caution.}
\label{tab:per_region}
\end{table}

Most Asian regions show elevated entropy relative to the Western baseline of 0.7437, with substantial variation across regions: Mongolia (0.8330) is most distant from ATU, while Malaysia (0.7248) is closest to the Western baseline. Myanmar/Shan, with only 9 examples, also shows relatively low entropy (0.7601) but is too underpowered for reliable inference. This regional variation is informative.

We form a preliminary hypothesis that geographic and cultural proximity to European traditions may correlate with ATU compatibility. Specifically, Malaysia's lower entropy is consistent with its position as a maritime trade hub; the Indian Ocean trade networks connected Malaysia to Indian, Arab, and Chinese merchant cultures, potentially facilitating narrative exchange and cultural contact that introduced patterns compatible with ATU categories.

In contrast, more geographically isolated regions like Mongolia (historically steppes cultures with distinct shamanic narrative traditions) show higher entropy and thus less compatibility with ATU. Japan, despite being a developed civilization, shows high entropy (0.8104), possibly reflecting either island geography limiting overland narrative exchange or distinct Buddhist narrative epistemologies. This pattern suggests that ATU compatibility may correlate with historical cultural contact rather than reflecting some fundamental universal narrative structure.

This regional breakdown is exploratory and necessarily tentative with current data. A rigorous test of the trade-route hypothesis would require independent quantification of historical cultural contact. Future work should examine regional patterns with multilingual data and explicit historical contact metrics.

\section{Analysis: Potential Confounds and Robustness}

Our finding---that Asian tales have lower ATU classifier confidence---is strong, but it is important to ask: \textit{Could this pattern result from confounds rather than genuine narrative incompatibility}? We carefully examine the major potential confounds and demonstrate that our results are robust.

\subsection{Confound 1: Text Length}

A critical observation: \textbf{Asian folktales are substantially longer than Western tales}. Our analysis shows:

\begin{itemize}
\item Western tales: mean 4,325 characters, median 2,843, max 59,583.
\item Asian tales: mean 8,329 characters, median 5,812, max 64,117.
\item Statistical test: Mann-Whitney U p-value $< 10^{-36}$, Cohen's $d = 0.621$ (large effect).
\end{itemize}

At first glance, this might seem problematic: if longer tales have richer linguistic features, they might produce higher classification confidence, which would work against finding lower Asian confidence. However, deeper analysis shows that this confound actually \textbf{strengthens rather than undermines our conclusion}:

\begin{enumerate}
\item \textbf{TF-IDF should benefit from length.} Longer texts produce more diverse TF-IDF token coverage. More tokens mean more features that could match in-distribution patterns, potentially increasing confidence. If Asian tales (longer) produce lower confidence \textit{despite} this length advantage, the true incompatibility is even stronger.

\item \textbf{Sentence embeddings truncate at 512 tokens.} Tales longer than 512 tokens (approximately 400 English words) are truncated. Asian tales average 8,329 characters, roughly equivalent to 1,800--2,200 words. Thus many Asian tales are substantially truncated, losing semantic information. This truncation actually \textit{harms} Asian confidence relative to Western tales, making our result conservative (understating the true gap).

\item \textbf{Result is conservative.} Both the TF-IDF and embedding arguments suggest that text length should favor Asian confidence (or at least not harm it). The fact that we observe lower Asian confidence \textit{despite} this advantage indicates our findings are conservative.
\end{enumerate}

\textbf{Conclusion:} The text length confound actually cuts in \textit{our favor}, making our result more conservative, not less.

\subsection{Confound 2: Translation Register and Era}

All texts are in English. Western folktales were collected and published primarily in the late 19th--early 20th centuries. Asian folktales were \textit{translated} by Western scholars of the same era (Skeat, Fleeson, Ozaki, etc.), who may have imposed European narrative conventions on Asian tales to make them more accessible to Western audiences.

\textbf{Critical insight:} If Victorian translators imposed European narrative conventions on Asian tales, this would make Asian tales \textit{more} ATU-compatible, not less. Our finding of lower Asian confidence \textit{despite} this bias suggests that genuine narrative incompatibility is even stronger than observed. The translation register bias thus also works in favor of our conclusion, making it conservative.

\subsection{Confound 3: Source Heterogeneity}

Western tales come from a single curated source (AFT); Asian tales come from approximately 15 different sources of varying provenance and quality. Could differences in source (rather than narrative tradition) drive our results?

\textbf{Evidence against source confound:}
\begin{enumerate}
\item \textbf{Multiple sources per region.} For Japan, we have tales from three independent sources (Ozaki, Nixon, Gutenberg curated). All show similar entropy (0.8104), suggesting the pattern is robust across sources.

\item \textbf{Per-region consistency.} All nine Asian regions show elevated entropy. No single source dominates the result. If a particular source were skewing results, we would expect regional variation; instead, the pattern is systematic.

\item \textbf{Methodologically sound.} We do not attempt to hide or control for source differences. Real-world corpora are heterogeneous; imposing artificial homogeneity would be misleading. Our confidence metrics are designed to be robust across sources.
\end{enumerate}

\textbf{Remaining limitation:} We cannot fully rule out source effects without a controlled study using identical tales translated multiple ways. This remains a useful direction for future work.

\subsection{Confound 4: Classifier Miscalibration (ECE = 0.20)}

We observed Expected Calibration Error of 0.20, indicating moderate miscalibration in the softmax probability scores. Could miscalibration be driving spurious results?

\textbf{Robustness to miscalibration:}

\begin{enumerate}
\item \textbf{Mann-Whitney U is rank-based.} The rank-based test is invariant to any monotonic transformation (e.g., temperature scaling or Platt sigmoid rescaling). Miscalibration preserves rank ordering, so the test is unaffected.

\item \textbf{Multi-metric agreement.} We report four independent metrics: softmax (probability-based), entropy (probability-based), margin (probability-based), and Mahalanobis (embedding-based). The Mahalanobis distance is computed entirely in embedding space and is unaffected by softmax calibration. That all four metrics show the same direction of effect (p < $3 \times 10^{-18}$) provides strong evidence that miscalibration is not driving the result.

\item \textbf{Non-parametric testing.} By using rank-based rather than parametric tests, we avoid assumptions about the distributional shape of confidence scores. This robustness is a key advantage of our statistical approach.
\end{enumerate}

\textbf{Limitation acknowledged:} ECE = 0.20 means that probability scores should not be interpreted as absolute calibrated probabilities. However, for our relative comparison (Western vs. Asian), the rank-based test is robust. We note this limitation and address it methodologically.

\subsection{Confound 5: CV Model vs. Final Model Asymmetry}

Western confidence comes from CV-fold models (80\% training data). Asian confidence comes from the final model (100\% training data). Could the final model's strength lead to inflated Asian confidence?

\textbf{Analysis:} More training data typically increases model confidence on \textit{all} data, including both in-distribution and out-of-distribution examples. If the final model is stronger, we would expect it to be more confident on Asian tales. Yet we observe it is \textit{less} confident. This suggests:

\begin{enumerate}
\item The true ATU incompatibility of Asian tales is large enough to overcome the model-strength advantage.
\item Our result is conservative: the gap would be even larger if we used equally-strong models for both Western and Asian inference.
\end{enumerate}

We can bound the strength difference: on the Western held-out test set (Level 1, 70\% training), we achieved 69.74\% accuracy, very close to the CV mean of 69.43\%. The improvement from 80\% to 100\% training data is small (approximately 0.3 percentage points). Thus the final model is not dramatically stronger, and this asymmetry is unlikely to substantially affect our conclusions.

\section{Discussion}

\subsection{Summary of Findings}

Our analysis provides strong quantitative evidence that the Aarne-Thompson-Uther Index, while highly effective for classifying European folktales (69.4\% cross-validated accuracy), does not generalize well to East and Southeast Asian folktales. The classifier's confidence drops significantly when applied to Asian tales across all four measured confidence metrics (normalized entropy rising from 0.7437 to 0.8025, $p < 3 \times 10^{-27}$, Cohen's $d \approx -0.56$). This gap persists when controlling for text length (which actually favors Asian confidence), translation register bias (which would make Asian tales artificially more ATU-compatible), source heterogeneity, classifier miscalibration, and the asymmetry between CV-fold and final models. In fact, each of these potential confounds, when analyzed, suggests that our findings are \textit{conservative}---the true narrative incompatibility between ATU and Asian traditions may be even larger.

The per-region breakdown reveals intriguing patterns: regions with historical cultural contact via trade routes (Tibet, Malaysia) show higher ATU compatibility than more geographically isolated regions (Mongolia, Japan). This suggests that ATU compatibility may correlate with historical exposure to European and Indo-European narrative traditions, rather than reflecting fundamental universal narrative structure.

\subsection{Implications for Folklore Scholarship}

\textbf{1. ATU is not culturally universal.} The ATU Index has become so foundational to folklore studies that scholars often apply it unreflectively to tales from any tradition. Our work provides empirical evidence that this universalist assumption is incorrect. The ATU system captures European narrative patterns exceptionally well but is less effective for East Asian traditions.

\textbf{2. The ATU Index embeds Western cultural conventions.} The categorical structure of ATU---emphasizing magical objects, supernatural helpers, enchanted spouses, and linear problem-solution narrative arcs---reflects European fairy-tale conventions shaped by Christian and Indo-European worldviews. East Asian narratives, shaped by Buddhist and Confucian philosophies, emphasize different narrative structures: cyclical spiritual transformation, moral enlightenment through ordeal, karma and fate, recursive and fractal narratives (as in the frame-story structure of stories-within-stories in classical Eastern literature).

\textbf{3. Implications for ``universal'' narrative claims.} In literary criticism and cognitive science, scholars have long proposed supposedly universal narrative structures (Campbell's monomyth, Propp's morphology, story grammar approaches). Our findings suggest that claims of universalism should be qualified: these may describe narrative structures common to Western/Indo-European traditions rather than genuinely universal human narrative architecture.

\textbf{4. Toward culturally-grounded narrative classification.} Rather than imposing ATU categories on non-Western tales, folklore scholarship would benefit from developing region-appropriate classification systems. East Asian traditions have their own tale-type systems (e.g., Japanese oral tradition scholars have developed context-specific classifications). Integrating computational methods with region-specific narrative expertise could produce more culturally grounded and nuanced classifications.

\subsection{Limitations}

Our study has several important limitations that should be acknowledged:

\begin{enumerate}
\item \textbf{No Europe-to-Europe transfer baseline.} To definitively show that lower Asian confidence reflects cultural mismatch rather than ordinary domain shift, we should train on one set of European regions and test on held-out European regions. Without this control, we cannot conclusively rule out that the pattern is simply what cross-source European domain shift looks like.

\item \textbf{Single classifier architecture.} We use LinearSVC with TF-IDF+embeddings. Other classifiers (deep neural networks, tree ensembles, transformer-based models) might show different confidence patterns. We claim robustness to confounds (text length, miscalibration), not robustness to architectural choice. Evaluation with multiple independent classifiers would strengthen generalizability claims.

\item \textbf{English translations only.} All texts are English translations, predominantly from Victorian and Edwardian-era sources. Translation artifacts, translator biases, and era-specific conventions affect linguistic patterns. Access to original-language texts and modern translations would enable cleaner tests of whether transferability failure is linguistic or conceptual.

\item \textbf{Limited fine-grained analysis.} Our approach operates at the tale level using distributional features (embeddings, TF-IDF). A more fine-grained analysis of narrative motifs, character types, plot structures, and thematic elements would provide insight into \textit{which} aspects of ATU categories transfer poorly versus well.

\item \textbf{Small sample sizes for some regions.} Myanmar/Shan (n=9) is substantially underpowered for reliable inference and should be interpreted with strong caveats. Regional patterns should ideally be validated on larger corpora before being treated as genuine phenomena.

\item \textbf{No human validation.} We lack ground-truth annotation of "ease of ATU fit" from folklore experts. Even a small scholar-annotated subset (50--100 tales) scoring how naturally each tale fits ATU categories would validate the bridge between model confidence and category fit.

\item \textbf{Snapshot study on fixed data.} This is a single evaluation on fixed corpora with a fixed label grouping. Replication on newly-collected Asian tale corpora, evaluation with alternative ATU groupings (finer or coarser), and evaluation with newly-collected European tales would test whether the findings generalize beyond this specific data slice.

\item \textbf{Transferability vs. categorical incompatibility.} This work demonstrates imperfect transferability of an ATU-trained classifier to Asian tales under this operationalization. It does not prove that ATU is fundamentally incompatible with Asian narrative traditions, only that the learned Western-to-Asian mapping is imperfect under our specific design choices.

\end{enumerate}

\subsection{Future Directions}

Several promising directions emerge from this work:

\begin{enumerate}
\item \textbf{Multilingual embeddings and original-language texts.} Using multilingual sentence transformers and analyzing tales in original languages (Chinese, Japanese, Korean, Tibetan, Mongolian, etc.) would allow us to assess whether translation register is a confound and whether narrative incompatibility is fundamental or translation-artifact.

\item \textbf{Fine-grained motif analysis.} Rather than relying on surface-level text patterns, future work could analyze tales at the level of narrative motifs. Existing motif ontologies (ATU itself, Stith Thompson's motif index) could be expanded and computationally analyzed to identify which motifs are characteristic of each tradition.

\item \textbf{Cross-cultural classifier training.} Train classifiers on mixed European + Asian corpora and observe whether the model learns balanced representations or defaults to European patterns. Analyze classifier representations to identify narrative dimensions along which cultures differ.

\item \textbf{Development of region-specific indices.} Work with regional folklore scholars to develop ATU-like classification systems for East Asian traditions. Integrate computational methods to systematize these region-specific indices.

\item \textbf{Quantification of cultural contact.} Using historical data on trade routes, diplomatic exchanges, and cultural contact, quantify the relationship between cultural proximity and narrative structure similarity. Validate the hypothesis that cultural contact correlates with narrative overlap.

\end{enumerate}

\section{Conclusion}

We set out to audit the transferability of the Aarne-Thompson-Uther Index across cultural domains. A text classifier trained on 1,518 Western European folktales labeled with coarse ATU categories achieves 69.4\% cross-validated accuracy on held-out Western tales, indicating that the label structure is learnable in-distribution. When applied to 453 East and Southeast Asian folktales from 8 regions, the same classifier shows significantly lower confidence across all four measured metrics (normalized entropy rising from 0.7437 to 0.8025, p < $3 \times 10^{-27}$, Cohen's $d \approx -0.56$).

\textbf{What this result does show:} ATU label structure, operationalized through text embeddings and TF-IDF features, transfers imperfectly from Western European tales to the evaluated Asian corpora. This imperfect transferability is robust across multiple confidence metrics, multiple regions, and multiple confound analyses.

\textbf{What this result does not show:} We have not proven that East and Southeast Asian folktales are categorically outside the ATU system in absolute terms, nor have we ruled out all forms of domain shift. We have tested one classifier family on translated texts under one operationalization of the 14-class ATU grouping. Other architectures (e.g., transformer-based), original-language texts, or finer-grained ATU categories might behave differently. Our findings characterize imperfect transferability under these specific conditions, not universal incompatibility.

\textbf{Implications for evaluation and benchmarks:} This work illustrates a broader lesson for machine learning evaluation: taxonomies that are convenient for one domain or historical context may encode domain-specific assumptions that limit their transportability. When benchmark systems or label schemes are developed in one population or cultural context, their transferability to other contexts should be empirically tested. This is not a problem unique to folklore; it applies to evaluation systems across domains. Audit methods like ours---combining in-distribution validation with out-of-distribution confidence measurement---offer practical tools for identifying taxonomy-transfer failures before they propagate into downstream applications.

Future work should extend this framework to other historically situated taxonomies and explore whether similar patterns of imperfect transferability arise in other domains. The development of region-appropriate narrative classification systems, grounded in cultural expertise and validated with computational methods, represents one promising path forward for more culturally attentive evaluation practices.

\begin{ack}

We thank the Archive of Finnish Tales (AFT) for providing the curated Western folktale corpus with ATU annotations. We acknowledge Project Gutenberg for digitizing and preserving the Victorian-era folklore translations that form the backbone of our Asian folktale collection. We are grateful to the Aarne-Thompson-Uther scholarly committee for establishing the foundational taxonomy that motivated this cross-cultural investigation. Finally, we thank the many folklore collectors, translators, and archivists whose labor over centuries made this analysis possible. This work received no specific grant or funding.

\end{ack}

\medskip

\section*{References}

{
\small

\begin{thebibliography}{99}

\bibitem[Aarne \& Thompson(1961)]{aarne1961}
Aarne, A., \& Thompson, S. (1961). \textit{The Types of the Folktale: A classification and bibliography} (2nd rev.). FF Communications No. 184.

\bibitem[Cohen(1988)]{cohen1988}
Cohen, J. (1988). \textit{Statistical power analysis for the behavioral sciences} (2nd ed.). Lawrence Erlbaum Associates.

\bibitem[Cortes \& Vapnik(1995)]{cortes1995}
Cortes, C., \& Vapnik, V. (1995). Support-vector networks. \textit{Machine Learning}, 20(3), 273--297.

\bibitem[Guo et al.(2017)]{guo2017}
Guo, C., Pleiss, G., Sun, Y., \& Weinberger, K. Q. (2017). On calibration of modern neural networks. \textit{International Conference on Machine Learning} (pp. 1321--1330). PMLR.

\bibitem[Kohavi(1995)]{kohavi1995}
Kohavi, R. (1995). A study of cross-validation and bootstrap for accuracy estimation and model selection. \textit{IJCAI} (Vol. 14, pp. 1137--1145).

\bibitem[Ledoit \& Wolf(2004)]{ledoit2004}
Ledoit, O., \& Wolf, M. (2004). A well-conditioned estimator for large-dimensional covariance matrices. \textit{Journal of Multivariate Analysis}, 88(2), 365--411.

\bibitem[Lee et al.(2018)]{lee2018}
Lee, K., Lee, K., Lee, H., \& Shin, J. (2018). A simple unified framework for detecting out-of-distribution samples and adversarial attacks. \textit{Advances in Neural Information Processing Systems} (pp. 7167--7177).

\bibitem[Mann \& Whitney(1947)]{mann1947}
Mann, H. B., \& Whitney, D. R. (1947). On a test of whether one of two random variables is stochastically larger than the other. \textit{Annals of Mathematical Statistics}, 18(1), 50--60.

\bibitem[Naeini et al.(2015)]{naeini2015}
Naeini, M. P., Cooper, G. F., \& Hauskrecht, M. (2015). Obtaining well calibrated probabilities using bayesian binning. \textit{AAAI} (Vol. 29, No. 1).

\bibitem[Platt(1999)]{platt1999}
Platt, J. (1999). Probabilistic outputs for support vector machines and comparisons to regularized likelihood methods. \textit{Advances in Large Margin Classifiers}, 10(3), 61--74.

\bibitem[Reimers \& Gurevych(2019)]{reimers2019}
Reimers, N., \& Gurevych, I. (2019). Sentence-BERT: Sentence embeddings using Siamese BERT-networks. \textit{Proceedings of the 2019 Conference on Empirical Methods in Natural Language Processing} (pp. 3982--3992). ACL.

\bibitem[Salton \& Buckley(1988)]{salton1988}
Salton, G., \& Buckley, C. (1988). Term-weighting approaches in automatic text retrieval. \textit{Information Processing \& Management}, 24(5), 513--523.

\bibitem[Thompson(1946)]{thompson1946}
Thompson, S. (1946). \textit{The folktale}. Dryden Press.

\bibitem[Uther(2004)]{uther2004}
Uther, H. J. (2004). \textit{The types of international folktales: A classification and bibliography} (3 vols.). FF Communications No. 284--286.

\bibitem[Zadrozny \& Elkan(2002)]{zadrozny2002}
Zadrozny, B., \& Elkan, C. (2002). Transforming classifier scores into accurate multiclass probability estimates. \textit{Proceedings of the Eighth ACM SIGKDD International Conference on Knowledge Discovery and Data Mining} (pp. 694--699).

\end{thebibliography}

}

\end{document}
